%==============================================================================
% CONFIG: STYLING (HEADING, PENOMORAN, HEADER-FOOTER, TOC, CAPTION, LIST)
%==============================================================================

%------------------------------------------------------------------------------
% PENOMORAN: Bab (I, II, ...), subbab 1.1., sub-subbab 1.1.1.
%------------------------------------------------------------------------------
\renewcommand{\thechapter}{\Roman{chapter}}
\renewcommand{\thesection}{\arabic{chapter}.\arabic{section}}
\renewcommand{\thesubsection}{\arabic{chapter}.\arabic{section}.\arabic{subsection}}
\renewcommand{\thefigure}{\arabic{chapter}.\arabic{figure}}
\renewcommand{\thetable}{\arabic{chapter}.\arabic{table}}
\setcounter{secnumdepth}{2}
\setcounter{tocdepth}{2}

%------------------------------------------------------------------------------
% HEADING FORMAT
%   Heading 1 (Chapter): rata tengah, tebal, "Bab I" lalu nama bab di baris baru
%   Heading 2 (Section): tebal, nomor 1.1.
%   Heading 3 (Subsection): tebal, nomor 1.1.1.
%------------------------------------------------------------------------------
\titleformat{\chapter}[display]
  {\normalfont\bfseries\centering\spasisetengah}
  {Bab \thechapter}{0pt}{}
\titlespacing*{\chapter}{0pt}{0pt}{12pt}

\titleformat{\section}[hang]
  {\normalfont\bfseries\spasisetengah}
  {\makebox[1.27cm][l]{\thesection.}}{0pt}{}
\titlespacing*{\section}{0pt}{3pt}{3pt}

\titleformat{\subsection}[hang]
  {\normalfont\bfseries\spasisetengah}
  {\makebox[1.27cm][l]{\thesubsection.}}{0pt}{}
\titlespacing*{\subsection}{0pt}{12pt}{3pt}

% Judul halaman khusus (misal: lembar pengesahan, daftar lain-lain)
\newcommand{\judulhal}[1]{%
  \begingroup\centering\bfseries #1\par\endgroup\vspace{3pt}}

%------------------------------------------------------------------------------
% NOMOR HALAMAN (PAGINATION)
%   Halaman pertama bab/lampiran : bottom center
%   Halaman berikutnya           : top right
%------------------------------------------------------------------------------
\fancypagestyle{halamanisi}{%
  \fancyhf{}%
  \renewcommand{\headrulewidth}{0pt}%
  \renewcommand{\footrulewidth}{0pt}%
  \fancyhead[R]{\thepage}%
}

\fancypagestyle{plain}{%
  \fancyhf{}%
  \renewcommand{\headrulewidth}{0pt}%
  \renewcommand{\footrulewidth}{0pt}%
  \fancyfoot[C]{\thepage}%
}

\pagestyle{halamanisi}

%------------------------------------------------------------------------------
% CAPTION: Tabel di atas tabel, Gambar di bawah gambar, format <bab>.<urutan>
%------------------------------------------------------------------------------
\captionsetup{labelsep=period,justification=centering,singlelinecheck=true,
              font=normalsize,labelfont=normalfont}
\captionsetup[table]{position=above,skip=6pt}
\captionsetup[figure]{position=below,skip=6pt}

%------------------------------------------------------------------------------
% DAFTAR ISI (TOC) DENGAN TOCLOFT
%------------------------------------------------------------------------------
\renewcommand{\cfttoctitlefont}{\bfseries}
\renewcommand{\cftaftertoctitle}{}
\renewcommand{\cftchapfont}{\bfseries}
\renewcommand{\cftchappagefont}{\bfseries}
\renewcommand{\cftchappresnum}{Bab }
\renewcommand{\cftdotsep}{0.6}   % Titik leader rapat
\renewcommand{\cftchapleader}{\bfseries\cftdotfill{\cftdotsep}}
\setlength{\cftchapnumwidth}{3.4em}
\setlength{\cftbeforechapskip}{3pt}
\setlength{\cftbeforesecskip}{3pt}
\setlength{\cftbeforesubsecskip}{3pt}
\setlength{\cftsecindent}{0.635cm}
\setlength{\cftsecnumwidth}{1.0cm}
\renewcommand{\cftsecaftersnum}{.}
\setlength{\cftsubsecindent}{1.27cm}
\setlength{\cftsubsecnumwidth}{1.3cm}
\renewcommand{\cftsubsecaftersnum}{.}

%------------------------------------------------------------------------------
% LIST / PENOMORAN ITEM
%------------------------------------------------------------------------------
\setlist{nosep}
\setlist[enumerate]{leftmargin=1.27cm,labelwidth=0.635cm,labelsep=0pt,align=left,
                    itemsep=0pt,topsep=0pt,parsep=0pt}
